\documentclass[journal]{IEEEtran}
\usepackage{amsmath,amssymb,amsfonts}
\usepackage{graphicx,booktabs,array,cite}
\usepackage{url}
\usepackage[hidelinks]{hyperref}
\hypersetup{pdftitle={Channel Estimation for Rydberg Receivers Using Learned Hankel Filtering},pdfauthor={Abdullah Haatim},pdfsubject={Research draft; simulation study},pdfkeywords={Rydberg receivers, channel estimation, Hankel filtering, self-supervised learning}}
\graphicspath{{figures/}}
\interdisplaylinepenalty=2500
\renewcommand{\topfraction}{0.92}
\renewcommand{\textfraction}{0.06}
\renewcommand{\floatpagefraction}{0.85}
\renewcommand{\dbltopfraction}{0.92}
\newcommand{\CN}{\mathcal{CN}}
\newcommand{\Hank}{\mathcal{H}}
\DeclareMathOperator{\sigmoid}{sigmoid}
\DeclareMathOperator{\softplus}{softplus}
\DeclareMathOperator{\tr}{tr}
\hyphenation{Ryd-berg mea-sure-ments}

\begin{document}
\title{Channel Estimation for Rydberg Receivers\\Using Learned Hankel Filtering}
\author{Abdullah~Haatim%
\thanks{A. Haatim is with the Department of Electrical and Electronics
Engineering, Birla Institute of Technology and Science, Pilani, India
(e-mail: f20220519@pilani.bits-pilani.ac.in).}
\thanks{Research draft, September 23, 2026. Results are simulation only.}}
\markboth{RESEARCH DRAFT, SEPTEMBER 2026}{HAATIM: CHANNEL ESTIMATION USING LEARNED HANKEL FILTERING}
\maketitle

\begin{abstract}
Rydberg receivers use excited atoms to sense radio signals. In a magnitude-only
receiver, the measured signal strength does not directly reveal the wireless
channel. Recovering the channel is difficult because the measurements contain
noise and do not directly provide phase. Channel structure can help, but a
fixed structural rule may also remove useful information. This paper studies
a small learning-based estimator that adjusts the amount of Hankel filtering
using the current measurements. It alternates between fitting the receiver
observations and filtering the channel estimate. Training uses a withheld
part of the observations, without true channels as targets. In simulations,
the method gives lower average channel error and shorter processing time than
the tested classical Hankel estimator. Further tests show that using both
channel-structure information and noise-related information improves average
accuracy. The benefit is not uniform: the classical method remains better
with a weak reference signal and in some channels with few paths. The results
support learned Hankel filtering as a useful option for the tested receiver
model, while showing the need to state its operating limits clearly.
\end{abstract}

\begin{IEEEkeywords}
Channel estimation, Hankel filtering, magnitude measurements, Rydberg
receivers, self-supervised learning.
\end{IEEEkeywords}

\section{Introduction}
\IEEEPARstart{R}{ydberg} atoms are atoms in highly excited states that can
respond to radio-frequency electric fields. This makes them useful for
studying new ways to receive wireless signals. In the receiver model
considered here, signals from several users are combined with a known
reference signal. The receiver measures the magnitude of that combined
signal, rather than its full complex value. The wireless channel describes
how each user's signal changes in strength and phase before reaching the
receiver. Estimating it is an important step in recovering the transmitted
information. Both model-based and learning-based methods have been developed
for this problem, including URformer~\cite{xu2025,xiao2026}.

Signals arriving through a small number of paths often form a repeated
pattern across a regular receiver array. A Hankel matrix arranges the channel
samples so that this pattern can be used for estimation. However, a channel
does not always follow a simple structural model closely. Filtering too
strongly can remove real channel details, while filtering too weakly can
leave noise in the estimate. Trying many possible structural settings also
takes time. This motivates a small model that learns how much filtering to
apply from the available measurements, without requiring the true channel
during training.

The contributions of this paper are:
\begin{enumerate}
\item A 188-parameter estimator that combines Rician measurement updates
with smooth Hankel filtering. A shared controller adjusts the filtering
strength using the estimated channel spectrum, residual and known noise level.
\item An application of held-out-measurement training to this estimator.
True channels are not training targets or checkpoint-selection scores.
The contribution is the receiver-specific estimator, not a new principle
of self-supervised learning.
\item Paired simulation comparisons of accuracy, processing time and
simpler versions of the method, including fresh-test replication and
input-removal tests. Weak-reference and few-path failures are reported.
\end{enumerate}

\begin{figure*}[!t]
\centering
\includegraphics[width=\textwidth]{fig1_method.pdf}
\caption{Estimator overview. The shared controller uses the current
estimate and known receiver information to choose filtering strength and
mixing. During training, the nine withheld pilots are used only to score
predictions; they do not enter initialization or the ten estimation stages.}
\label{fig:method}
\end{figure*}

\section{Related Work}
Learning without clean target signals is already used in reconstruction.
One approach uses part of the measurements to form an estimate and another
part to train the model~\cite{yaman2020}. Self-supervised learning has also
been used to remove noise from Rydberg sensor signals~\cite{liu2026}.
Huang et al.\ learn Hankel-factor updates from fully sampled signal
targets~\cite{huang2021}. Luo et al.\ combine self-supervision and
structured low-rank reconstruction for linearly sampled complex MRI
data~\cite{luo2023}. DeepTensor represents low-rank factors by networks
fitted to the observed object~\cite{saragadam2022}. Rician-likelihood
training for magnitude data is also established~\cite{parker2023}.
Thus neither self-supervision, learned Hankel processing nor the Rician
loss is new on its own.

Our focus is a particular estimation task: recovering a complex multiuser
channel from noisy magnitudes with a known reference. We combine the
receiver's noise model with a small network that controls smooth Hankel
filtering. Our companion manuscript studies classical Hankel estimation and
rank selection~\cite{haatimcompanion}. Here, the emphasis is on learning
filtering strength without channel labels and measuring accuracy and
processing cost. Unlike factor-generating networks, our controller
sets two filtering controls; a fixed ten-stage estimator performs
the channel inversion. Unlike methods that restore measured complex
samples directly, it fits noisy magnitudes through a known reference.
These are design differences, not a claim of superiority over those
methods. Self-supervised low-rank channel estimation in XL-MIMO is also
listed in the CCAI 2026 program~\cite{ccai2026}. Its full method was not
available for a detailed comparison. We make no first-ever claim and
do not claim to reproduce the published URformer results.

\section{Receiver and Channel Model}
Consider a receiver array with $N$ elements and $K$ users. Each user sends
$P$ known training symbols, called pilots. These pilots allow the receiver
to estimate the channel before detecting unknown data. The observations are
\begin{equation}
 Z=|GS+B+W|,\qquad W_{np}\sim\CN(0,v),
 \label{eq:model}
\end{equation}
where $G\in\mathbb C^{N\times K}$ is the unknown channel,
$S\in\mathbb C^{K\times P}$ contains the pilots, and
$B\in\mathbb C^{N\times P}$ is the known reference. Noise samples are
independent. The magnitude acts on each entry, so $Z$ is nonnegative and
does not directly measure phase. The noise variance $v$ is the total
complex variance; each real component has variance $v/2$.

The reference makes the measured magnitudes depend on the relative phases
of the user signals. It does not guarantee that every channel can be
recovered accurately. Reference calibration and noise variance are assumed
known, except in tests that deliberately supply an incorrect noise variance.
Receiver conversion gain is absorbed into the model's scaling. These are
simulation assumptions, not claims about all physical Rydberg receivers.

The main channel uses a half-wavelength uniform linear array:
\begin{equation}
 g_{nk}=\sum_{\ell=1}^{L_k}\alpha_{\ell k}
 e^{-j\pi n\sin\theta_{\ell k}},\quad n=0,\ldots,N-1.
 \label{eq:channel}
\end{equation}
Each path has gain $\alpha_{\ell k}$ and arrival angle $\theta_{\ell k}$.
Angles are uniform on $[-\pi/2,\pi/2]$, and independent gains follow
$\CN(0,1/L_k)$, giving unit average power per user. Pilot entries follow
$\CN(0,1)$. The reference has a known spatial response and constant
commanded waveform across the pilot observations.

SNR measures total average user power relative to noise. RSR measures
reference power relative to one user's average power. For the unit-power
users used here,
\begin{equation}
 \mathrm{SNR}=10\log_{10}(K/v),\quad
 \mathrm{RSR}=10\log_{10}|\alpha_b|^2,
\end{equation}
where $\alpha_b$ is the reference amplitude. With three users, an RSR of
12~dB is about 7.23~dB relative to their combined average signal power.
These definitions must remain separate when interpreting the results.

\section{Proposed Estimator}
\subsection{Updating the Estimate from Measurements}
The estimator starts from a channel estimate formed from the available
observations, pilots and reference. It repeats a measurement update and a
Hankel-filtering step ten times, as illustrated in Fig.~\ref{fig:method}.

The measurement update uses the Rician distribution, which describes the
magnitude of a complex signal corrupted by the assumed Gaussian noise.
For signal magnitude $a=|\mu|$, its density is
\begin{equation}
 p(z\mid a,v)=\frac{2z}{v}
 \exp\left(-\frac{z^2+a^2}{v}\right)I_0(2za/v).
\end{equation}
Here $I_0$ is a modified Bessel function. This model accounts for noise
acting before the magnitude is taken. It does not replace the magnitude
observations with an additive Gaussian-noise approximation.

Let $M=GS+B$ and $\kappa=2Z|M|/v$, with entrywise operations in
\begin{equation}
 Q=Z\frac{I_1(\kappa)}{I_0(\kappa)}\frac{M}{|M|}.
\end{equation}
$I_1$ is the first-order modified Bessel function. The continuous limit is
used at $M=0$. The expectation-maximization (EM) update is
\begin{align}
 G_{\rm EM}&=(Q-B)S^H(SS^H)^{-1},\\
 G_{\rm dc}&=G+\gamma_t(G_{\rm EM}-G),
 \label{eq:dc}
\end{align}
where $\gamma_t=2\sigmoid(a_t)$ is a learned step size at stage $t$.
The inverse is evaluated by a linear solve. The input pilot matrix must
have full row rank. The implementation requires at least $2K$ input pilots;
this is a numerical design rule, not a proof of unique channel recovery.

\subsection{Smooth Hankel Filtering}
Each user's channel vector in $G_{\rm dc}$ is arranged into a balanced
Hankel matrix $H=\Hank(g_{{\rm dc},k})$, whose entries are constant along
antidiagonals. At $N=32$, its size is $17\times16$. Its singular values
describe the strengths of different matrix components. Their collection
is called the singular-value spectrum.

A hard rank rule keeps a chosen number of components and discards the
rest. We instead reduce each singular value smoothly:
\begin{equation}
 \widehat\sigma_i=\frac{\sigma_i^3}{\sigma_i^2+\tau^2}.
 \label{eq:shrink}
\end{equation}
Larger $\tau$ gives stronger filtering. Small singular values are reduced
by a larger fraction than large ones. This does not mean that every small
component is noise: some may contain real channel information.

Writing $q=\lceil N/2\rceil$ and
$s_H=\|H\|_F/\sqrt q$, the controller output $u$ sets
\begin{equation}
 \tau=s_H\bigl[0.001+\softplus(u)\bigr].
\end{equation}
The floor makes the normalized solve nonsingular. The equivalent matrix
operation is
\begin{equation}
 \widehat H=H-\tau^2H(H^HH+\tau^2I)^{-1}.
\end{equation}
The code normalizes by $s_H$ before solving and uses a small scale guard
at zero. Gradients pass through this solve without an identity-backward
approximation. Antidiagonal averaging, denoted by $\Hank^\dagger$, converts
the filtered matrix back to a channel vector:
\begin{equation}
 g_k^{t+1}=g_{{\rm dc},k}+\eta_{t,k}
 \left(\Hank^\dagger\widehat H_k-g_{{\rm dc},k}\right).
\end{equation}
The learned mixing factor $\eta_{t,k}=\sigmoid(w)$ lies between zero and
one. Neither $\eta$ nor $\tau$ is an estimate of the true path count.
The complete learned iteration is not guaranteed to reduce channel error
at every stage or to converge.

\subsection{Choosing the Filtering Strength}
A shared network with eight inputs, 16 hidden units and two outputs
chooses $u$ and $w$. It uses a SiLU activation, $x\sigmoid(x)$. The same
network is used for every user and stage. It has 178 parameters; ten
learned step sizes bring the total to 188.

The first three inputs describe the Hankel spectrum. If
$p_i=\sigma_i/\sum_j\sigma_j$, they are
\begin{equation}
 \frac{\exp(-\sum_i p_i\log p_i)}{q},\qquad p_1,\qquad
 \sum_{i=q-\lceil q/2\rceil+1}^{q}p_i.
\end{equation}
These describe how spread out the components are, the largest component's
share, and the smaller components' combined share.

Two inputs describe noise and the difference between measured and predicted
magnitudes, called the measurement residual:
\begin{align}
 f_4&=\log\left(1+\frac{\operatorname{mean}[(Z-|G_{\rm dc}S+B|)^2]}{v}\right),\\
 f_{5,k}&=\log\left(1+\frac{v\tr[(SS^H)^{-1}]}{K\widehat p_k}\right),
 \quad \widehat p_k=\frac{\|g_{{\rm dc},k}\|^2}{N}.
\end{align}
The remaining inputs are $P/(2K)$, $\log q$ and $(t+1)/10$, with
$t=0,\ldots,9$. Small denominator guards prevent division by zero.
The model does not receive true channels, true path counts or channel-family
labels. Its outputs change with the observations, but its learned weights
remain fixed during use.

\subsection{Training Without True Channels}
Each training example has 30 pilot columns. A random split gives 21 columns
to the estimator and withholds nine. The estimate is used to predict the
withheld magnitudes. Training minimizes their mean negative log likelihood,
computed stably as
\begin{equation}
 \ell(z,\mu,v)=\frac{(z-|\mu|)^2}{v}
 -\log I_{0e}(2z|\mu|/v),
\end{equation}
where $I_{0e}$ is the exponentially scaled Bessel function. An observation-only
term is omitted because it is the same for the models being compared.
Small-field series evaluate the continuous limits stably. The loss fits
the amplitude parameter $|\mu|$ of the Rician density, not its mean;
these are different in noise.

This is self-supervised learning: the observations provide training
information without a true-channel target. Splits change each epoch, which
is one pass through the training examples. Separate validation examples and
a fixed split choose the saved model. At deployment, all available pilots
are used, with no extra measurements and no further weight training. This
changes the input pilot count from its training value, so performance under
the change is tested rather than assumed.

\section{Simulation Study}
\subsection{Setup and Error Measure}
Table~\ref{tab:setup} lists the main settings. Half the training examples
have random counts of three to seven paths per user. The rest are spread
over fixed counts of 2, 4, 8, 12 and 16. Validation gives equal weight to
126 SNR, RSR and path-setting combinations. The simulated condition labels,
including path settings, are used only to set these validation weights;
they are not estimator inputs. Model selection uses held-out likelihood,
not channel error.

\begin{table}[!t]
\caption{Main Simulation and Training Settings}
\label{tab:setup}
\centering\small
\begin{tabular}{@{}ll@{}}
\toprule
Setting & Value\\\midrule
Array elements / users / pilots & $32/3/30$\\
Training / validation cases & $16{,}384/2{,}048$\\
SNR values (dB) & $-10,-5,0,5,10,15,20$\\
Training RSR values (dB) & $5,12,20$\\
Training seeds & 3, sharing training cases\\
Estimator stages / parameters & $10/188$\\
Training epochs / batch size & $20/32$\\
Optimizer / learning rate & Adam / $0.001$\\
Input / withheld pilots & $21/9$\\
Training / evaluation precision & Single / double\\
\bottomrule
\end{tabular}
\end{table}

The main study trains 18 models: two update choices, three filtering
choices and three training seeds. Seeds change initial weights, batch
order and pilot splits, but share the same training cases. EM remains the
primary update because it was selected by validation likelihood before
test generation. The 18-model training campaign took about 3.67 hours of
elapsed time with overlapping workers, excluding data preparation and
preflight. This is a campaign cost, not the training time of one model.

The main evaluation covers 32 conditions with 400 paired cases per
condition. Two noise-input tests reuse cases from another condition, giving
12,000 distinct cases. Methods receive the same channel, pilots and noise
within each case. The main classical comparator is HS-GS, the
Hankel-structured iterative estimator with rank search from the companion
study~\cite{haatimcompanion}. For each case, HS-GS fits ranks $1,\ldots,q$
on 70\% of the pilots using 25 EM updates and four alternating rank/Hankel
projection sweeps per update. It selects one rank using squared magnitude
residuals on the remaining pilots, then refits using all pilots for 100
updates. The rank search consumes no extra observations and is included
in timing. Its selection loss differs from our Rician training loss.
Other controls include fixed-iteration EM,
Rician and amplitude-gradient methods, hard Hankel filtering and learned
versions with global or no filtering. The global control learns a relative
threshold and mixing factor per stage; its absolute threshold still scales
with $s_H$. It has 30 parameters. The no-filtering model has ten. The
hard-Hankel Rician-gradient control uses rank ten, selected on development
likelihood, and one projection sweep per update. Constant-output controls
replace the adaptive network's outputs with validation averages; they
are interventions, not separately trained models.

Channel accuracy is measured by normalized mean squared error (NMSE):
\begin{equation}
 \mathrm{NMSE}_{\rm dB}=10\log_{10}
 \frac{\sum_i\|\widehat G_i-G_i\|_F^2}{\sum_i\|G_i\|_F^2}.
\end{equation}
Lower values are better. We average the per-seed dB scores, not the
predicted channels. The main average gives equal weight to the seven
SNR settings. A positive gain means that the proposed method has lower
error than the comparator. Severe errors have per-case linear NMSE above one.

Uncertainty intervals use 4,000 paired resamples of whole test cases.
Main-study intervals are pointwise 95\% intervals, not simultaneous
guarantees across every tested condition. They describe test-case
variation for the saved models, not the full uncertainty from new training
datasets. Stopping-controlled classical methods have a 1,000-update cap;
reaching that cap is not treated as convergence.

\begin{figure}[!t]
\centering
\includegraphics[width=\columnwidth]{fig2_snr.pdf}
\caption{Main SNR comparison at $N=32$, $K=3$, $P=30$ and RSR 12~dB,
with random three-to-seven-path channels. Lower NMSE is better. Lines show
mean per-seed scores; shaded bands are conditional paired-world 95\%
intervals, not training-dataset uncertainty.}
\label{fig:snr}
\end{figure}

\subsection{Main Accuracy and Processing Time}
\begin{table}[!t]
\caption{All Main-Study Comparators}
\label{tab:all}
\centering\footnotesize
\setlength{\tabcolsep}{3pt}
\begin{tabular}{@{}lrrr@{}}
\toprule
Method & NMSE & Time & Severe\\
 & (dB) & (ms) & (\%)\\\midrule
\csname @@input\endcsname baseline_rows.tex
\bottomrule
\end{tabular}
\vspace{3pt}
\parbox{\columnwidth}{\footnotesize NMSE and severe-error frequency are
equal-weight averages across the seven primary SNR conditions. Time is
the median at SNR 5~dB, RSR 12~dB, not across the SNR sweep. Learned
scores average three seeds; classical cases are not counted three times.
All numbers describe the tested implementations.}
\end{table}

Table~\ref{tab:all} includes every main-study comparator. HS-GS is the
strongest tested classical method on the declared SNR average. The
secondary adaptive gradient-descent (GD) core has a slightly lower test
NMSE than the primary EM core. We retain EM because core selection was
made from validation likelihood, not from this test ranking. The
proposed method is therefore not the best row on every metric.

Figure~\ref{fig:snr} compares the proposed method with HS-GS and simpler
learned controls. Its average gain over HS-GS is 1.675~dB, with an interval
of $[1.618,1.730]$~dB. The three seed-specific gains range from 1.533 to
1.787~dB. The improvement is not equally strong at every SNR. At 20~dB,
the gain is 0.139~dB and its interval, $[-0.017,0.294]$~dB, includes zero.
That condition therefore does not establish a clear improvement.

\begin{figure}[!t]
\centering
\includegraphics[width=\columnwidth]{fig4_runtime.pdf}
\caption{Median processing time at SNR 5~dB and RSR 12~dB. Timing includes
initialization and rank search where required. One-thread Apple M2 CPU,
double precision, one case at a time, 100 cases and three repeats.
The asterisk denotes a stopping-controlled Rician gradient solver. These
are implementation-specific timings, not hardware-independent speed limits.}
\label{fig:time}
\end{figure}

Figure~\ref{fig:time} shows processing times measured under the same timing
protocol. The proposed estimator takes a median of 8.24~ms per case,
compared with 525.62~ms for HS-GS, a measured speedup of about 64 times.
This does not establish the same speed difference on other hardware or
against differently optimized software. Offline training is a separate
cost and is not included in inference time.
The learned implementations use PyTorch, while the classical ones use
NumPy/SciPy. The no-filtering learned model is faster still, but has
higher average error. The primary method's severe-error frequency is
0.45\%, compared with 14.25\% for HS-GS across the primary SNR cells.
These are descriptive rates for this mixture, not guarantees of
reliability for other operating conditions.

\begin{figure*}[!t]
\centering
\includegraphics[width=\textwidth]{fig3_limits.pdf}
\caption{Limits of the main method at SNR 5~dB. (a) Reference strength
with random three-to-seven-path channels. (b) Fixed path counts at RSR
12~dB. Positive gain favors the proposed estimator; negative gain favors
HS-GS. Error bars are pointwise paired-world 95\% intervals conditional
on the trained models. The path-count sweep overlaps training support.}
\label{fig:limits}
\end{figure*}

\subsection{Operating Limits}
Figure~\ref{fig:limits} includes settings where HS-GS remains better.
The proposed method loses 1.617~dB at RSR 0~dB, 1.104~dB with two paths,
and 0.203~dB with three paths. These are important limits. A model chosen
for good average performance need not be best for a very simple channel
or a weak reference. Using the correct noise distribution does not
guarantee that the estimation problem is easy to solve.

Separate transfer tests give gains for the tested clustered, ten-path
and unequal-power channels, and for supplied noise variances that are
wrong by $\pm3$~dB. Results also change with array size and user count;
not every condition has a clear advantage. Some path counts occur in
training, so the sweep is not evidence of success on wholly unfamiliar
channel types. These tests do not establish that the method resolves
failures observed in earlier experiments with different protocols.
The alternative-channel tests use four clusters of ten rays within
$\pm5^\circ$, ten independent equal-power paths, or ten paths with
geometrically decreasing powers. Each alternative is rescaled to the
corresponding base channel's total energy. They are one-condition stress
tests at SNR 5~dB and RSR 12~dB, not full channel-family SNR sweeps.

\subsection{Further Comparisons}
A separate follow-up trains 15 models and uses two fresh test sets, each
with seven main SNR settings and two difficult settings of 200 cases.
Six primary comparisons are adjusted for multiple testing within each
set. Table~\ref{tab:checks} reports the average gains. The independently
generated training datasets each use one fixed training seed, so they
provide limited replication rather than a full study of dataset and
initialization variation.

\begin{table}[!t]
\caption{Follow-up Average Gains (dB)}
\label{tab:checks}
\centering\small
\begin{tabular}{@{}p{.70\columnwidth}rr@{}}
\toprule
Comparison & Set 1 & Set 2\\\midrule
New-data adaptive vs. global filtering & 2.252 & 2.185\\
New-data adaptive vs. no filtering & 4.250 & 4.191\\
New-data adaptive vs. HS-GS & 1.939 & 1.837\\
Original adaptive vs. supervised recipe & 1.015 & 1.012\\
Original adaptive vs. final-only filtering & 0.697 & 0.694\\
Fixed-open mix vs. original adaptive & 0.064 & 0.066\\
\bottomrule
\end{tabular}
\end{table}

The supervised comparison uses the same architecture, input pilots and
training budget, but uses true channels as targets. Its training objective
is mean linear channel NMSE. The label-free model has about a 1~dB
advantage under the main averaging rule. However, selecting the supervised
checkpoint by held-out likelihood reduces this to about 0.73~dB. Pooling
linear errors across SNRs reduces it to about 0.36~dB. Supervision is
better in some settings, including the two-path test. It also transfers
less consistently from 21 input pilots to 30 in a validation replay.
We therefore do not conclude that self-supervision is generally better
than supervision.

Filtering throughout the estimator improves average accuracy by about
0.70~dB over applying filtering only at the end. It also costs more: a
separate timing sample gives about 4.65~ms and 2.05~ms, respectively.
These timings are not combined with the main study's measurements to
form a different speedup claim.

The learned mixing factor is usually close to one. Fixing it at one gives
only a small average gain and worsens two of three training seeds on
the two-path cases. The tested fixed-mix model still stores 188 parameters.
We keep the original method and do not describe the mixing factor as a
reliable decision to switch filtering on or off.

\begin{figure*}[!t]
\centering
\includegraphics[width=\textwidth]{fig5_features.pdf}
\caption{Smaller-budget controller-input study. (a) Average gain of the
full controller over versions with restricted inputs. Intervals adjust
for three comparisons within each test set and are conditional on the
trained models. (b) Full versus noise/residual-only inputs by SNR, showing
the loss at 5~dB. The dashed line marks the predeclared 0.25~dB protection
limit. Both fresh test sets are shown; each contains 128 cases per condition.}
\label{fig:features}
\end{figure*}

\subsection{Which Controller Inputs Help?}
A smaller study compares the full controller with spectrum-only,
noise/residual-only and size/stage-only versions. Size and stage inputs
are retained in every version. All versions use the same measurement
update and Hankel filter. Each is trained from scratch on the same
4,096 development cases, with three seeds and 12 epochs. Two fresh
test sets contain 2,304 cases in total: seven main SNR settings and two
difficult settings, with 128 cases per condition in each set.

Figure~\ref{fig:features}(a) shows average gains of 1.22--1.25~dB over
spectrum-only inputs, 0.57--0.58~dB over noise/residual-only inputs, and
2.21--2.23~dB over size/stage-only inputs. Gains are positive for all
three seeds in both sets. Intervals adjusted for the three comparisons
remain above zero. Final and validation-selected checkpoints give the
same scores in this study.

The full controller is not better everywhere. In
Fig.~\ref{fig:features}(b), it is about 0.39--0.40~dB worse than the
noise/residual-only version at 5~dB SNR. It therefore fails the predeclared
requirement that no main condition should worsen by more than 0.25~dB
against that version. In the separate two-path cases, it loses about
0.82--0.93~dB to the same version.

The results support combining both types of input to improve the average
under this training setup. They do not show that both are necessary in
every channel. The reduced budget also does not replace the main study.
All models store 188 parameters, but masked input weights have no effect.
Removing noise or spectral inputs from the controller does not remove
all related information: the measurement update still uses noise variance,
and the filter still acts on the estimated channel.

\section{Discussion and Conclusion}
The useful result is not that Hankel filtering always helps. It is that
a small model can learn how strongly to use this structure, without
true-channel training targets, and provide a useful average accuracy and
processing-time result in the tested setting. The input-removal study
supports using both channel-structure and noise-related information,
while its losses show why an average alone is insufficient.
The practical use case is repeated estimation under conditions similar
to the simulated training population, where offline training can be
reused. This is an accuracy--processing-cost result for the implemented
comparators, not a universal replacement for classical estimation.

All results are simulations. Reference calibration is assumed known,
and no physical receiver measurements are used. Three seeds sharing
training cases are not three independent datasets. We have not proved
unique channel recovery in every setting or convergence of the full
learned estimator. The comparisons also do not establish superiority
over published URformer performance.
Noise-dependent regularized classical controls and broader independent
training/pilot replications would be needed for a stronger robustness
claim. We do not infer the cause of a loss solely from its association
with SNR, path count or reference strength.

Within these limits, learned Hankel filtering offers a compact,
label-free option for magnitude-only channel estimation. Its gains and
its weak-reference and few-path failures should be considered together
when deciding whether to use it.

\appendix[Initialization and Numerical Checks]
For array element $n$, form $A_n\in\mathbb C^{P\times(K+1)}$ with row
$[s_p^T,\,b_{np}]$. Let $u_n$ be the dominant eigenvector of
$A_n^H\operatorname{diag}(z_n)A_n$ and set
\begin{equation}
 c_n=\frac{\sum_p z_{np}|(A_nu_n)_p|}
 {\sum_p |(A_nu_n)_p|^2}.
\end{equation}
The first $K$ components of
$c_nu_n\exp[-j\arg(u_{n,K+1})]$ give the initial channel row. During
training, this uses only the 21 input pilots. This initializer is part
of the shared estimation pipeline, not a new contribution.

Numerical checks cover the Rician likelihood and derivatives, complex
noise variance, finite gradients, repeated singular values, pilot-split
isolation, single/double-precision agreement and deterministic checkpoint
recovery. The latest input study passed its saved-output audit for all
432 result files and 24 checkpoint replays. Such checks verify the
implementation; they do not turn the failed protection criterion into
a successful result.

One finite-input SVD failure occurred in the classical HS-GS array-size
test ($N=64$). The same matrix was recomputed using SciPy's alternative
SVD driver; the recorded relative reconstruction error was below
$3\times10^{-15}$. The trial was retained, and its inputs, equations and
iteration budget were unchanged.

\section*{Reproducibility}
The accompanying reproducibility files include source code, frozen
configurations, model checkpoints, compact per-case errors, timing records
and source hashes. Recorded seeds specify the simulation datasets.
Full observation arrays and full predicted channels are omitted from
the compact package; this omission is listed explicitly. A small saved
observation subset supports checkpoint replay. The numerical environment
is Python 3.12 with NumPy 2.4.6, SciPy 1.17.1 and PyTorch 2.13.0.

\newpage
\section*{Acknowledgment}
OpenAI Codex~\cite{codex2026} assisted with drafting and revising all
sections, mathematical exposition, experiment and analysis code, and
preparation of Figs.~1--5. Numerical plots were generated from saved
simulation outputs. The author is responsible for the scientific claims,
references and final submitted content.

\bibliographystyle{IEEEtran}
\bibliography{references}
\end{document}
